\section{5. Theoretical Analysis}
To isolate the mechanical conditions that govern consensus failure and stability transitions under status-based penalization, this section analyzes a minimal network topology\cite{degroot1974reaching}. By evaluating the exact difference equations of a three-agent path graph, the dynamics are characterized across three analytical stages\cite{degroot1974reaching}: first, resolving the assimilation problem under pure DeGroot averaging ($\gamma_{stub} = 0$)\cite{degroot1974reaching}; second, deriving the step-by-step state trajectories and stationary fixed point under partial stubbornness ($\gamma_{stub} = 0.2$)\cite{degroot1974reaching}; and third, formalizing the piecewise affine geometry, border-collision bifurcations, and ergodicity properties that separate fixed dissensus from sustained limit cycles\cite{acemoglu2013opinion,hajnal1958weak}.

\subsection{The Minimal Network Setup}
Consider a directed path graph of three agents, $V = {1, 2, 3}$, configured in a linear communication chain, $1 \leftrightarrow 2 \leftrightarrow 3$, where agent 2 occupies the central structural position\cite{degroot1974reaching}.

The network is initialized with private scalar estimates\cite{degroot1974reaching}: $$x_1(0) = 60, \quad x_2(0) = 50, \quad x_3(0) = 40$$ The initial unweighted network mean is $\mu(0) = (60 + 50 + 40)/3 = 50$\cite{degroot1974reaching}. Agent 1 is an upper outlier whose estimate exceeds the group mean, designating it as a "tall poppy" subject to penalization\cite{degroot1974reaching}. The system is parameterized with maximal penalty magnitude $P = 1.0$, step threshold $\theta = 0.5$, and an influence step size $\epsilon = 0.5$\cite{degroot1974reaching}.

The baseline interaction protocol assigns every agent an endogenous self-weight of $1 - \epsilon = 0.5$\cite{degroot1974reaching}. The remaining weight pool $\epsilon = 0.5$ is distributed equally among all unpenalized incoming neighbors\cite{degroot1974reaching}. Using the convention that $\tilde{w}{ij}(t)$ represents the weight listener $i$ assigns to speaker $j$, the non-homogeneous affine update rule for agent $i$ is\cite{degroot1974reaching}: $$x_i(t+1) = \gamma{stub} x_i(0) + (1 - \gamma_{stub}) \sum_{j=1}^3 \tilde{w}_{ij}(t) x_j(t)$$

Under maximal penalty ($P = 1.0$), whenever an agent's standardized deviation satisfies $D_i(t) > 0.5$, its outgoing influence factor is zeroed: $f_i(t) = 0$\cite{degroot1974reaching,acemoglu2013opinion}. When penalized, listener 2 sets $\tilde{w}{21}(t) = 0$ and allocates its entire neighbor weight pool $\epsilon = 0.5$ to agent 3\cite{degroot1974reaching}. Conversely, when agent 1 is unpenalized ($f_1(t) = 1$), listener 2 splits its neighbor pool equally between agents 1 and 3, setting $\tilde{w}{21}(t) = 0.25$ and $\tilde{w}_{23}(t) = 0.25$\cite{acemoglu2013opinion}.

\subsection{The Assimilation Problem (\gamma_{stub} = 0)}
To establish why cognitive anchoring is an absolute mathematical requirement for non-convergent behavior, consider the system under pure social averaging without stubbornness ($\gamma_{stub} = 0$)\cite{degroot1974reaching}: $$x(t+1) = \tilde{W}(t) x(t)$$

At $t = 0$, the network statistics are $\mu(0) = 50$ and population standard deviation $\sigma(0) = \sqrt{\frac{1}{3}[(60-50)^2 + (50-50)^2 + (40-50)^2]} = \sqrt{200/3} \approx 8.165$\cite{acemoglu2013opinion}. The standardized deviation of agent 1 evaluates to\cite{acemoglu2013opinion}: $$D_1(0) = \frac{60 - 50}{\sqrt{200/3}} = \sqrt{\frac{3}{2}} \approx 1.225 > 0.5$$ Because $D_1(0) > 0.5$ and $P = 1.0$, agent 1 is fully penalized ($f_1(0) = 0$)\cite{acemoglu2013opinion}. Agents 2 and 3 sit at or below the mean ($D_2(0) = 0$, $D_3(0) = -\sqrt{3/2}$), retaining full throughput ($f_2(0) = f_3(0) = 1.0$)\cite{acemoglu2013opinion}.

The resulting transition matrix $\tilde{W}_A$ at $t = 0$ is\cite{acemoglu2013opinion}: $$\tilde{W}_A = \begin{bmatrix} 0.5 & 0.5 & 0.0 \ 0.0 & 0.5 & 0.5 \ 0.0 & 0.5 & 0.5 \end{bmatrix}$$

Evaluating the state update from $t = 0$ to $t = 1$\cite{acemoglu2013opinion}: $$x_1(1) = 0.5(60) + 0.5(50) = 55.0$$ $$x_2(1) = 0.5(50) + 0.5(40) = 45.0$$ $$x_3(1) = 0.5(50) + 0.5(40) = 45.0$$

At $t = 1$, agents 2 and 3 achieve identical states: $x_2(1) = x_3(1) = 45.0$\cite{acemoglu2013opinion}. The network mean drops to $\mu(1) = (55 + 45 + 45)/3 = 145/3 \approx 48.333$\cite{acemoglu2013opinion}. The standard deviation contracts to $\sigma(1) = 10\sqrt{2}/3 \approx 4.714$\cite{acemoglu2013opinion}. Evaluating agent 1's standardized status deviation at $t = 1$\cite{acemoglu2013opinion}: $$D_1(1) = \frac{55 - 145/3}{10\sqrt{2}/3} = \frac{20/3}{10\sqrt{2}/3} = \sqrt{2} \approx 1.414 > 0.5$$ Because $D_1(1) > 0.5$, the penalty remains fully active ($f_1(1) = 0$), and the transition matrix remains $\tilde{W}(1) = \tilde{W}_A$\cite{acemoglu2013opinion}.

For any time step $t \ge 1$ where $x_2(t) = x_3(t) = 45.0$ and $x_1(t) > 45.0$, the subsystem ${2, 3}$ remains static\cite{acemoglu2013opinion}: $$x_2(t+1) = 0.5(45) + 0.5(45) = 45.0$$ $$x_3(t+1) = 0.5(45) + 0.5(45) = 45.0$$ Agent 1 evolves according to the linear recurrence\cite{acemoglu2013opinion}: $$x_1(t+1) = 0.5 x_1(t) + 0.5(45.0) = 0.5 x_1(t) + 22.5$$

Solving this linear non-homogeneous difference equation with initial condition $x_1(1) = 55.0$ yields the exact closed-form trajectory for all $t \ge 1$\cite{acemoglu2013opinion}: $$x_1(t) = 45.0 + 10.0 \cdot (0.5)^{t-1} = 45.0 + 20.0 \cdot (0.5)^t$$ The unweighted collective network mean evolves as\cite{acemoglu2013opinion}: $$\mu(t) = \frac{x_1(t) + 45.0 + 45.0}{3} = 45.0 + \frac{20}{3} \cdot (0.5)^t$$

Calculating the standardized status deviation along this trajectory for any $t \ge 1$\cite{acemoglu2013opinion}: $$x_1(t) - \mu(t) = \frac{40}{3} \cdot (0.5)^t$$ $$\sigma(t) = \sqrt{\frac{1}{3} \left[ \left( \frac{40}{3} \cdot (0.5)^t \right)^2 + 2 \left( -\frac{20}{3} \cdot (0.5)^t \right)^2 \right]} = \frac{20\sqrt{2}}{3} \cdot (0.5)^t$$ $$D_1(t) = \frac{x_1(t) - \mu(t)}{\sigma(t)} = \frac{\frac{40}{3} \cdot (0.5)^t}{\frac{20\sqrt{2}}{3} \cdot (0.5)^t} = \frac{40}{20\sqrt{2}} = \sqrt{2} \approx 1.4142$$

Equation (7) reveals an invariance property: for all finite time steps $t \ge 1$, the standardized status deviation of agent 1 is identically constant at $D_1(t) = \sqrt{2}$\cite{acemoglu2013opinion}. Because $\sqrt{2} > 0.5$, the penalty factor remains locked at $f_1(t) = 0$ for all finite $t$\cite{acemoglu2013opinion}. Agent 1 never catches the falling group mean at any finite time; the state gap contracts asymptotically as $t \to \infty$\cite{acemoglu2013opinion}: $$\lim_{t \to \infty} x(t) = [45.0, ; 45.0, ; 45.0]^T = 45.0 \cdot \mathbf{1}$$

Under pure DeGroot averaging, status penalization destroys the network's accurate upper estimate without cycling\cite{degroot1974reaching}. Agent 1 is assimilated into the unpenalized subgroup consensus at $45.0$, producing a terminal consensus error of $|\mu(\infty) - \mu(0)| = |45 - 50| = 5.0$\cite{degroot1974reaching,acemoglu2013opinion}. This establishes that time-varying edge suppression alone cannot generate sustained limit cycles; non-convergent behavior requires an exogenous anchor holding the outlier in place\cite{degroot1974reaching}.

\subsection{Step-by-Step Transition Logic (\gamma_{stub} = 0.2)}
Introducing partial stubbornness at $\gamma_{stub} = 0.2$ anchors 20% of each agent's state update to its initial private estimate, while exposing the remaining 80% to social averaging\cite{degroot1974reaching}. The difference equations become\cite{degroot1974reaching}: $$x_1(t+1) = 0.2(60) + 0.8 \sum_{j=1}^3 \tilde{w}{1j}(t) x_j(t) = 12.0 + 0.8 \sum{j=1}^3 \tilde{w}{1j}(t) x_j(t)$$ $$x_2(t+1) = 0.2(50) + 0.8 \sum{j=1}^3 \tilde{w}{2j}(t) x_j(t) = 10.0 + 0.8 \sum{j=1}^3 \tilde{w}{2j}(t) x_j(t)$$ $$x_3(t+1) = 0.2(40) + 0.8 \sum{j=1}^3 \tilde{w}{3j}(t) x_j(t) = 8.0 + 0.8 \sum{j=1}^3 \tilde{w}_{3j}(t) x_j(t)$$

Round 1 ($t = 0 \to t = 1$)

At $t = 0$, $D_1(0) \approx 1.225 > 0.5$, triggering maximum suppression ($f_1(0) = 0$)\cite{degroot1974reaching}. The influence matrix is $\tilde{W}(0) = \tilde{W}_A$ from Equation (3)\cite{acemoglu2013opinion}. Evaluating the updates\cite{degroot1974reaching}: $$x_1(1) = 12.0 + 0.8 [0.5(60) + 0.5(50)] = 12.0 + 0.8(55.0) = 12.0 + 44.0 = 56.0$$ $$x_2(1) = 10.0 + 0.8 [0.5(50) + 0.5(40)] = 10.0 + 0.8(45.0) = 10.0 + 36.0 = 46.0$$ $$x_3(1) = 8.0 + 0.8 [0.5(40) + 0.5(50)] = 8.0 + 0.8(45.0) = 8.0 + 36.0 = 44.0$$ The network mean evaluates to\cite{degroot1974reaching}: $$\mu(1) = \frac{56.0 + 46.0 + 44.0}{3} = \frac{146.0}{3} \approx 48.667$$ Because $x_1(1) = 56.0 > 48.667$, $D_1(1) > 0.5$, holding the penalty active into Round 2\cite{degroot1974reaching}.

Round 2 ($t = 1 \to t = 2$)

Using the states at $t = 1$ under matrix $\tilde{W}_A$\cite{degroot1974reaching}: $$x_1(2) = 12.0 + 0.8 [0.5(56.0) + 0.5(46.0)] = 12.0 + 0.8(51.0) = 12.0 + 40.8 = 52.8$$ $$x_2(2) = 10.0 + 0.8 [0.5(46.0) + 0.5(44.0)] = 10.0 + 0.8(45.0) = 10.0 + 36.0 = 46.0$$ $$x_3(2) = 8.0 + 0.8 [0.5(44.0) + 0.5(46.0)] = 8.0 + 0.8(45.0) = 8.0 + 36.0 = 44.0$$ The network mean evaluates to\cite{degroot1974reaching}: $$\mu(2) = \frac{52.8 + 46.0 + 44.0}{3} = \frac{142.8}{3} = 47.60$$ Agent 1's estimate of $52.8$ exceeds $47.60$, maintaining the penalty into Round 3\cite{degroot1974reaching}.

Round 3 ($t = 2 \to t = 3$)

Advancing to $t = 3$\cite{degroot1974reaching}: $$x_1(3) = 12.0 + 0.8 [0.5(52.8) + 0.5(46.0)] = 12.0 + 0.8(49.4) = 12.0 + 39.52 = 51.52$$ $$x_2(3) = 10.0 + 0.8 [0.5(46.0) + 0.5(44.0)] = 10.0 + 0.8(45.0) = 10.0 + 36.0 = 46.0$$ $$x_3(3) = 8.0 + 0.8 [0.5(44.0) + 0.5(46.0)] = 8.0 + 0.8(45.0) = 8.0 + 36.0 = 44.0$$ The network mean evaluates to\cite{degroot1974reaching}: $$\mu(3) = \frac{51.52 + 46.0 + 44.0}{3} = \frac{141.52}{3} \approx 47.173$$

Notice that agents 2 and 3 reach their stationary fixed points on the very first round ($x_2^* = 46.0, x_3^* = 44.0$) because their convex neighbor combination evaluates to $0.5(46) + 0.5(44) = 45.0$\cite{acemoglu2013opinion}. Agent 1's descent slows monotonically: $60.0 \to 56.0 \to 52.8 \to 51.52$\cite{degroot1974reaching}.

\subsection{The Mechanism of Non-Assimilation}
The state trace demonstrates why assimilation fails when $\gamma_{stub} > 0$\cite{degroot1974reaching}. While the group mean drifts downward ($50.0 \to 48.67 \to 47.60 \to 47.17$), the gap between agent 1 and the mean narrows ($10.0 \to 7.33 \to 5.20 \to 4.35$) but remains strictly bounded away from zero\cite{degroot1974reaching,acemoglu2013opinion}.

The stubborn anchor acts as an irreducible mathematical floor\cite{degroot1974reaching}. During the penalty phase, agent 1's update is\cite{degroot1974reaching}: $$x_1(t+1) = \gamma_{stub} x_1(0) + (1 - \gamma_{stub}) [0.5 x_1(t) + 0.5 x_2(t)]$$ Because $x_2(t) = 46.0$ for all $t \ge 1$, agent 1 converges asymptotically to the stationary limit $x_1^$ satisfying $x_1^ = 12.0 + 0.4 x_1^* + 0.4(46.0)$\cite{degroot1974reaching,acemoglu2013opinion}: $$(1 - 0.4) x_1^* = 12.0 + 18.4 = 30.4 \implies x_1^* = \frac{30.4}{0.6} = \frac{152}{3} \approx 50.667$$

At this fixed point, the terminal network mean is\cite{degroot1974reaching,acemoglu2013opinion}: $$\mu^* = \frac{50.667 + 46.0 + 44.0}{3} = \frac{422}{9} \approx 46.889$$ The terminal population standard deviation evaluates to\cite{acemoglu2013opinion}: $$\sigma^* = \sqrt{\frac{1}{3} \left[ \left( \frac{152}{3} - \frac{422}{9} \right)^2 + \left( 46 - \frac{422}{9} \right)^2 + \left( 44 - \frac{422}{9} \right)^2 \right]} = \frac{\sqrt{632}}{9} \approx 2.7933$$ Agent 1's terminal standardized status deviation evaluates to\cite{acemoglu2013opinion}: $$D_1^* = \frac{x_1^* - \mu^}{\sigma^} = \frac{34/9}{\sqrt{632}/9} = \frac{34}{\sqrt{632}} \approx 1.3524$$

Because $D_1^* \approx 1.3524 > 0.5$, the penalty condition remains permanently satisfied ($f_1^* = 0$)\cite{acemoglu2013opinion}. The network topology never switches back to the unpenalized state\cite{acemoglu2013opinion}. The network freezes into a stationary equilibrium of persistent disagreement, formally establishing Regime 1B (Fixed Dissensus)\cite{acemoglu2013opinion}.

\subsection{Invariant Partitions, Fixed Dissensus, and Border-Collision Bifurcations}
The three-agent system constitutes a piecewise affine discrete dynamical system governed by two distinct topological modes\cite{acemoglu2013opinion,hajnal1958weak}: $$x(t+1) = \begin{cases} (1 - \gamma_{stub}) \tilde{W}A x(t) + \gamma{stub} x(0), & \text{if } D_1(x(t)) > \theta \quad (\text{System } A: \text{Penalty Active}) \ (1 - \gamma_{stub}) \tilde{W}B x(t) + \gamma{stub} x(0), & \text{if } D_1(x(t)) \le \theta \quad (\text{System } B: \text{Unpenalized}) \end{cases}$$

The transition matrices for the two regimes are\cite{acemoglu2013opinion}: $$\tilde{W}_A = \begin{bmatrix} 0.5 & 0.5 & 0.0 \ 0.0 & 0.5 & 0.5 \ 0.0 & 0.5 & 0.5 \end{bmatrix}, \qquad \tilde{W}_B = \begin{bmatrix} 0.5 & 0.5 & 0.0 \ 0.25 & 0.5 & 0.25 \ 0.0 & 0.5 & 0.5 \end{bmatrix}$$

Each linear mode possesses an isolated virtual fixed point $x_M^* = (I - (1 - \gamma_{stub}) \tilde{W}M)^{-1} \gamma{stub} x(0)$ for $M \in {A, B}$\cite{acemoglu2013opinion,hajnal1958weak}. For $\gamma_{stub} = 0.5$ and initial vector $x(0) = [60, 50, 40]^T$, the virtual fixed points evaluate to\cite{acemoglu2013opinion}: $$x_A^* = \left( I - 0.5 \tilde{W}_A \right)^{-1} [30, 25, 20]^T = [55.833, ; 47.50, ; 42.50]^T$$ $$x_B^* = \left( I - 0.5 \tilde{W}_B \right)^{-1} [30, 25, 20]^T = [56.667, ; 50.00, ; 43.333]^T$$

Evaluating the standardized deviation of agent 1 at these virtual fixed points\cite{acemoglu2013opinion}: $$D_1(x_A^) \approx 1.3132, \qquad D_1(x_B^) = \sqrt{\frac{3}{2}} \approx 1.2247$$

For the tolerance threshold $\theta = 0.5$, both virtual fixed points lie on the same side of the switching boundary\cite{acemoglu2013opinion}: $$D_1(x_A^) > 0.5 \quad \text{and} \quad D_1(x_B^) > 0.5$$ Because $D_1(x_A^) > \theta$, the fixed point $x_A^$ resides in the interior of its own active domain $\mathcal{D}A = {x \in \mathbb{R}^3 : D_1(x) > \theta}$\cite{acemoglu2013opinion}. Because $|(1 - \gamma{stub}) \tilde{W}A|\infty = 1 - \gamma_{stub} < 1$, the operator is a strict contraction mapping\cite{acemoglu2013opinion,hajnal1958weak}. Any trajectory initialized within $\mathcal{D}_A$ converges monotonically to $x_A^*$ without ever encountering the switching boundary $\mathcal{S} = {x : D_1(x) = \theta}$\cite{acemoglu2013opinion}.

A self-sustaining limit cycle (Regime 3) requires a border-collision bifurcation\cite{acemoglu2013opinion,hajnal1958weak}. In switched affine systems, persistent periodic orbits emerge if and only if the switching surface $\mathcal{S}$ separates the state space such that each subsystem's virtual fixed point lies on the opposite side of the boundary\cite{acemoglu2013opinion}: $$D_1(x_A^) < \theta < D_1(x_B^)$$

When Equation (17) is satisfied, neither subsystem can settle at its fixed point\cite{acemoglu2013opinion}:

While in the penalty-active mode ($\mathcal{D}_A$), the trajectory contracts toward $x_A^*$, which pulls $D_1(x(t))$ downward below $\theta$, crossing the switching surface into $\mathcal{D}_B$\cite{degroot1974reaching,acemoglu2013opinion}.

Upon entering $\mathcal{D}B$, the penalty lifts, restoring weight $\tilde{w}{21} > 0$ (the restitution phase)\cite{degroot1974reaching,acemoglu2013opinion}.

In $\mathcal{D}_B$, the trajectory contracts toward $x_B^*$, which pulls $D_1(x(t))$ upward above $\theta$, forcing the trajectory back across $\mathcal{S}$ into $\mathcal{D}_A$\cite{degroot1974reaching,acemoglu2013opinion}.

In generalized multi-agent networks ($N \ge 20$), this grazing condition is driven endogenously: as unpenalized lower agents contract, the network standard deviation $\sigma(t)$ compresses, dynamically rescaling $D_i(t)$ and driving the system across the switching boundary\cite{acemoglu2013opinion,hajnal1958weak}.

\subsection{Derivation of the Self-Sustaining Cycling Condition}
To establish the parameter boundary where downward assimilation is arrested, we examine the force balance acting on agent 1 during the penalty phase\cite{degroot1974reaching,acemoglu2013opinion}. The single-step displacement of agent 1 decomposes into two competing forces\cite{acemoglu2013opinion}: $$x_1(t+1) - x_1(t) = \gamma_{stub} \big(x_1(0) - x_1(t)\big) - (1 - \gamma_{stub}) \epsilon \big(x_1(t) - x_2(t)\big)$$ Here, $\gamma_{stub} (x_1(0) - x_1(t)) \ge 0$ is the restoring force exerted by the cognitive anchor, while $(1 - \gamma_{stub}) \epsilon (x_1(t) - x_2(t)) \ge 0$ is the downward drag exerted by social averaging with agent 2\cite{acemoglu2013opinion}.

For agent 1 to halt downward assimilation and maintain a positive displacement relative to the drifting collective, the upward restoring force must counterbalance the downward social drag\cite{acemoglu2013opinion}: $$\gamma_{stub} \big(x_1(0) - x_1(t)\big) \ge (1 - \gamma_{stub}) \epsilon (1 - \tilde{w}{21}) \big(x_1(t) - x_2(t)\big)$$ where the factor $(1 - \tilde{w}{21})$ parameterizes the severity of outgoing suppression enforced by listener 2\cite{degroot1974reaching,acemoglu2013opinion}.

Under maximum suppression ($P = 1.0$), $\tilde{w}{21} = 0$\cite{degroot1974reaching}. Under the heuristic assumption that the displacement from the anchor balances the distance to the neighboring state ($x_1(0) - x_1(t) \approx x_1(t) - x_2(t)$), dividing both sides yields the classical force-balance condition\cite{degroot1974reaching,acemoglu2013opinion}: $$\gamma{stub} > (1 - \gamma_{stub}) \epsilon \implies \gamma_{stub} > \frac{\epsilon}{1 + \epsilon}$$

Substituting the baseline influence step size $\epsilon = 0.5$ yields the analytical threshold\cite{degroot1974reaching}: $$\gamma_{stub} > \frac{0.5}{1 + 0.5} = \frac{1}{3} \approx 0.333$$

In rigorous dynamical systems terms, the true threshold separating fixed dissensus from limit cycling is governed by the position of the stationary fixed point relative to the switching manifold $\mathcal{S}$\cite{acemoglu2013opinion,hajnal1958weak}. Solving the full fixed-point equation $x_A^(\gamma_{stub}) = (I - (1 - \gamma_{stub}) \tilde{W}A)^{-1} \gamma{stub} x(0)$ gives\cite{acemoglu2013opinion}: $$x_2^(\gamma_{stub}) = \frac{10 \gamma_{stub} + 36 (1 - \gamma_{stub})}{1 + 0.2 (1 - \gamma_{stub})}, \qquad x_3^(\gamma_{stub}) = \frac{8 \gamma_{stub} + 36 (1 - \gamma_{stub})}{1 + 0.2 (1 - \gamma_{stub})}$$ $$x_1^(\gamma_{stub}) = \frac{60 \gamma_{stub} + 0.4 (1 - \gamma_{stub}) x_2^*(\gamma_{stub})}{1 - 0.4 (1 - \gamma_{stub})}$$

A border-collision bifurcation occurs at the critical value $\gamma_{stub}^$ where the stationary status deviance intersects the tolerance threshold\cite{acemoglu2013opinion}: $$D_1\big(x_A^(\gamma_{stub}^)\big) = \frac{x_1^(\gamma_{stub}^) - \mu^(x_A^)}{\sigma^(x_A^)} = \theta$$ When $\gamma_{stub} > \gamma_{stub}^$, cognitive anchoring halts downward assimilation, preserving the outlier gap required to destabilize the fixed point and initiate sustained limit cycles\cite{degroot1974reaching,acemoglu2013opinion}.

\subsection{Implications for Ergodicity and the Hajnal Coefficient}
The dynamical properties of the switched affine system provide direct insight into the failure of matrix ergodicity\cite{degroot1974reaching}.

In the minimal network setup under maximum suppression, listener 2 severs incoming influence from agent 1 ($\tilde{w}_{21} = 0$)\cite{degroot1974reaching}. However, this column zeroing does not force the instantaneous Hajnal coefficient to $\delta(\tilde{W}_A) = 1.0$\cite{acemoglu2013opinion}. Evaluating $\delta(\tilde{W}_A)$ from Equation (3)\cite{acemoglu2013opinion}: $$\frac{1}{2} |(\tilde{W}A){1\cdot} - (\tilde{W}A){2\cdot}|_1 = \frac{1}{2} (|0.5 - 0| + |0.5 - 0.5| + |0 - 0.5|) = 0.5$$ $$\frac{1}{2} |(\tilde{W}A){1\cdot} - (\tilde{W}A){3\cdot}|_1 = \frac{1}{2} (|0.5 - 0| + |0.5 - 0.5| + |0 - 0.5|) = 0.5$$ $$\frac{1}{2} |(\tilde{W}A){2\cdot} - (\tilde{W}A){3\cdot}|_1 = 0.0$$ $$\delta(\tilde{W}_A) = \max(0.5, ; 0.5, ; 0.0) = 0.5 < 1.0$$

The Dobrushin scrambling coefficient is strictly positive\cite{acemoglu2013opinion}: $$\alpha(\tilde{W}_A) = 1 - \delta(\tilde{W}_A) = 0.5 > 0$$ Agent 2 maintains a self-weight of $0.5$ and broadcasts to both agents 1 and 3 with weight $0.5$, providing universal common support across all three rows\cite{acemoglu2013opinion}. The matrix $\tilde{W}_A$ remains strictly scrambling\cite{acemoglu2013opinion}.

The breakdown of consensus does not arise from an open-loop failure of matrix scrambling\cite{acemoglu2013opinion,hajnal1958weak}. Rather, it stems from the algebraic structure of the affine system in Equation (1)\cite{acemoglu2013opinion,hajnal1958weak}. In the augmented coordinate space $\tilde{x}(t) = [x(t)^T, 1]^T$, the augmented transition operator is\cite{acemoglu2013opinion}: $$\tilde{H}(t) = \begin{bmatrix} (1 - \gamma_{stub}) \tilde{W}(t) & \gamma_{stub} x(0) \ \mathbf{0}^T & 1 \end{bmatrix}$$

Because the $(N+1)$-th coordinate is an invariant absorbing state, the dominant eigenspace of $\tilde{H}(t)$ is tied to the private vector $x(0)$\cite{acemoglu2013opinion}. When the border-collision condition $D_1(x_A^) < \theta < D_1(x_B^)$ is met, the system alternates infinitely often between matrices $\tilde{H}A$ and $\tilde{H}B$\cite{degroot1974reaching,acemoglu2013opinion}: $$\lim{T \to \infty} \prod{\tau=0}^{T-1} \tilde{H}(T - 1 - \tau) \ne \mathbf{1} v^T$$ The backward matrix product fails to converge to a rank-one matrix\cite{degroot1974reaching,acemoglu2013opinion}. State trajectories remain trapped in a non-convergent, bounded limit cycle, precluding asymptotic consensus and formalizing ergodicity failure in status-suppressed networks\cite{degroot1974reaching,acemoglu2013opinion}.

